% =========================================================================
% appendix_tables.tex — Appendix B
% =========================================================================

\section{Supplementary Tables}
\label{app:tables}

% -----------------------------------------------------------------
\subsection{NIAH Recall by Context Length and Needle Depth}

\begin{table}[H]
\centering
\caption{NIAH recall (\%) by scheme, context length, and needle depth.
All values from \texttt{results/mc\_niah.json} and
\texttt{results/standard\_eval.json}.}
\label{tab:niah_full}
\small
\begin{tabular}{@{}lcccccccccccc@{}}
\toprule
& \multicolumn{3}{c}{\textbf{256}} & \multicolumn{3}{c}{\textbf{512}}
& \multicolumn{3}{c}{\textbf{768}} & \multicolumn{3}{c}{\textbf{1024}} \\
\cmidrule(lr){2-4}\cmidrule(lr){5-7}\cmidrule(lr){8-10}\cmidrule(lr){11-13}
\textbf{Scheme} & E & M & L & E & M & L & E & M & L & E & M & L \\
\midrule
Baseline FP16
  & 100 & 100 & 100 & 100 & 100 & 100 & 100 & 100 & 100 & 100 & 100 & 100 \\
TQ Flat
  & 100 & 100 & 100 & 90 & 100 & 80 & 100 & 100 & 100 & 100 & 90 & 100 \\
WHT+Asym
  & 100 & 100 & 100 & 100 & 100 & 100 & 100 & 90 & 100 & 100 & 90 & 100 \\
MC-TQ
  & 100 & 100 & 100 & 100 & 100 & 90 & 100 & 100 & 90 & 100 & 100 & 90 \\
\bottomrule
\end{tabular}
% Source: results/mc_niah.json — all values verbatim
\end{table}

% -----------------------------------------------------------------
\subsection{Distortion Bounds: MSE and Inner Product}

\begin{table}[H]
\centering
\caption{Empirical MSE and inner-product distortion versus Shannon lower
bounds at $\dhead = 64$ for TurboQuant (TQ), WHT + Lloyd-Max, and
Outlier-split schemes.
All values from \texttt{results/distortion\_bounds.json}.}
\label{tab:distortion_bounds}
\small
\begin{tabular}{@{}ccccccc@{}}
\toprule
\textbf{Bits} & \textbf{Shannon} & \textbf{TQ MSE} & \textbf{WHT MSE}
  & \textbf{$\frac{\text{TQ}}{\text{Shannon}}$}
  & \textbf{$\frac{\text{WHT}}{\text{Shannon}}$}
  & \textbf{Outlier MSE} \\
\midrule
2 & 0.0625 & 0.1147 & 0.1146 & 1.83 & 1.83 & 0.2741 \\
3 & 0.0156 & 0.0334 & 0.0334 & 2.14 & 2.14 & 0.0869 \\
4 & 0.0039 & 0.0091 & 0.0091 & 2.34 & 2.34 & 0.0252 \\
5 & 0.00098 & 0.0024 & 0.0024 & 2.47 & 2.46 & 0.0069 \\
\bottomrule
\end{tabular}
% Source: results/distortion_bounds.json — mse_results array
%   2-bit: shannon=0.0625, tq=0.1147, wht=0.1146, out=0.2741
%   3-bit: shannon=0.0156, tq=0.0334, wht=0.0334, out=0.0869
%   4-bit: shannon=0.0039, tq=0.0091, wht=0.0091, out=0.0252
%   5-bit: shannon=0.00098, tq=0.0024, wht=0.0024, out=0.0069
\end{table}

% -----------------------------------------------------------------
\subsection{MC Compressor Ablation}

\begin{table}[H]
\centering
\caption{Checkpoint versus independent segment-state compression
in Memory Caching.
The checkpoint variant accumulates quantization error across segments,
degrading both perplexity and NIAH recall.
All values from \texttt{results/mc\_compressor\_ablation.json}.}
\label{tab:ablation}
\small
\begin{tabular}{@{}lcccc@{}}
\toprule
\textbf{Strategy} & \textbf{Avg PPL} & \textbf{Speed (tok/s)}
  & \textbf{Cache (KB)} & \textbf{NIAH Avg (\%)} \\
\midrule
Checkpoint & 17.06 & 27.96 & 13{,}056 & 0.0 \\
Independent & 2.13 & 30.63 & 6{,}912 & 96.7 \\
\bottomrule
\end{tabular}
% Source: results/mc_compressor_ablation.json
%   checkpoint: avg_ppl=17.064, speed=27.957, cache=13056, niah all 0%
%   independent: avg_ppl=2.134, speed=30.628, cache=6912
%     niah: 11 of 12 = 100%, 768/Late = 80% → avg = 96.67%
\end{table}

% -----------------------------------------------------------------
\subsection{LongBench-lite Results}

\begin{table}[H]
\centering
\caption{LongBench-lite subtask scores (\%). Zero-values (represented by a mathematical dash ``---'') represent zero-shot formatting and parsing failures inherent to the base model's pre-training limitations under context truncation, rather than execution bugs in our quantization script. All values from \texttt{results/longbench\_lite.json}.}
\label{tab:longbench}
\small
\begin{tabular}{@{}lccc@{}}
\toprule
\textbf{Scheme} & \textbf{Passage QA} & \textbf{Summ.\ R-1} & \textbf{Code Compl.} \\
\midrule
Baseline FP16   & 25.0  & 11.99 & --- \\
WHT+Asym 4-bit  & ---   & ---   & 5.0 \\
Outlier 2.5-bit & 5.0   & ---   & 5.0 \\
Outlier 3.5-bit & ---   & ---   & 5.0 \\
\bottomrule
\end{tabular}
% Source: results/longbench_lite.json — all values verbatim
\end{table}

% -----------------------------------------------------------------
\subsection{FibQuant Comparison}

\begin{table}[H]
\centering
\caption{FibQuant versus TurboQuant MSE at $d = 64$, four-bit.
The joint 2D radial-angular codebook is inferior on both MSE distortion
and perplexity.
All values from \texttt{results/fibquant\_results.json}.}
\label{tab:fibquant}
\small
\begin{tabular}{@{}lcc@{}}
\toprule
\textbf{Scheme} & \textbf{MSE Distortion} & \textbf{Avg PPL} \\
\midrule
TurboQuant MSE & 0.00179 & 4.62 \\
FibQuant 2D & 0.00338 & 5.05 \\
\bottomrule
\end{tabular}
% Source: results/fibquant_results.json →
%   turboquant_mse: mse=0.00179, ppl=4.620
%   fibquant_2d: mse=0.00338, ppl=5.048
\end{table}

% -----------------------------------------------------------------
\subsection{Complexity-Recall Tradeoff}

\begin{table}[H]
\centering
\caption{Memory Caching complexity-recall tradeoff across segment sizes.
Recall remains 100\% for segment sizes 8--64 and 256, dropping to 90\%
only at segment size~128.
All values from \texttt{results/complexity\_recall.json}.}
\label{tab:complexity}
\small
\begin{tabular}{@{}cccc@{}}
\toprule
\textbf{Segment Size} & \textbf{Cost Fraction}
  & \textbf{Recall (\%)} & \textbf{Cache (KB)} \\
\midrule
8   & 0.125   & 100.0 & 15{,}036 \\
16  & 0.0625  & 100.0 & 15{,}576 \\
32  & 0.03125 & 100.0 & 16{,}656 \\
64  & 0.0156  & 100.0 & 18{,}816 \\
128 & 0.0078  &  90.0 & 23{,}136 \\
256 & 0.0039  & 100.0 & 31{,}776 \\
\bottomrule
\end{tabular}
% Source: results/complexity_recall.json — all values verbatim
\end{table}
